Managing a live equity portfolio raises three linked questions. An allocation rule answers what to hold, given
forecasts and limits on risk, concentration and trading; a rebalancing rule answers when the holdings should be
recomputed; and an evaluation asks whether the two rules together would have worked on markets other than the one
that was observed. Mean-variance optimization~\cite{markowitz1952} and its continuous-time
extensions~\cite{merton1969} give a principled answer to the first question, but optimized weights are sensitive to
errors in the estimated means and covariances~\cite{michaud1989, ledoitWolf2004}, and an equal-weight portfolio
often matches optimized portfolios out of sample~\cite{demiguelGarlappiUppal2009}. Transaction costs turn the second
question into a control problem. With proportional costs, the optimal policy leaves the holdings alone inside a
no-trade region and trades only at its boundary~\cite{magillConstantinides1976, davisNorman1990}, and multi-period
formulations solve a sequence of forecast-based optimizations over a receding horizon~\cite{herzogDondiGeering2007,
boydEtAl2017}. The third question is usually answered with a backtest on the historical record, which is a single
path of the market. Choosing among rules on one path inflates their measured performance~\cite{baileyEtAl2014,
harveyLiu2015}, and a short window may contain few of the market conditions a portfolio must survive.

These lines of work mostly treat the three questions one at a time, although a live portfolio answers them
together and the answers interact. The allocation rule decides which positions the rebalancing rule must track, and
the market's path between rebalances decides when the next trade comes and from which holdings it starts. An
allocation that is easy to interpret helps here, because a change in a holding can then be traced to a change in one
stated number. Utility functions from consumer theory offer such a rule when each stock's preference weight is a
score with a plain meaning~\cite{cobbDouglas1928, arrowEtAl1961, stone1954, geary1950}. Testing the combined system
beyond the historical path requires synthetic markets that reproduce the stylized facts of daily growth rates, in
particular heavy tails and volatility clustering~\cite{mandelbrot1963, cont2001}, while keeping the dependence
between stocks and the market~\cite{sharpe1963}. Per-asset hidden Markov generators with a jump mechanism, combined
through a common market factor with each asset's own draw rescaled so that the market term does not inflate the
asset's total variance, reproduce the heavy tails of individual stocks within a multi-asset
market~\cite{alswaidan2026variance}. Such a market makes it possible to run a complete system, with its calibration
repeated on each simulated history, on futures that did not occur, including futures with deeper declines than the historical
window contains.

In this study, we built a portfolio management system that answers the three questions together and tested it on
real and synthetic markets. A selection rule chose the stocks a portfolio could hold, by ranking a score across the
whole universe or within each sector, by a bandit within each sector, or at random as a control, and two short lists
chosen in advance by the author served as further selections. Each stock's score divided its expected growth rate
from a single-index fit to the market by its absolute beta raised to a power that was positive when the market had
recently been falling and negative when it had been rising, so that, among stocks with positive expected growth,
low-beta stocks scored higher after declines and high-beta stocks after rises, and a hyperbolic tangent kept the score bounded. Utility functions from
consumer theory, with the scores as preference weights, and three standard portfolios (minimum variance, tangency
and long-short) set the holdings under shared limits on the budget, each stock's weight, portfolio
volatility and trading cost. An optional market hedge capped the portfolio's beta, each time the allocation was
re-solved, more tightly when the market was weaker and added up to five stocks with zero or negative beta to those the portfolio could hold. A model-predictive controller forecast the value of the
holdings with a multi-asset geometric Brownian motion and re-solved the allocation from the current holdings when
its horizon had passed, when the value left a band around the forecast, when a drawdown limit was reached, or when
a scheduled pick changed the stocks. We
calibrated the system on a decade of United States equity prices and ran it on the following twenty months, and we
ran it again on simulated price paths of the same period (synthetic futures) built from per-asset hidden Markov
generators with jumps, including futures, selected from many simulated paths, with deeper market declines than the
historical window contained. On real prices, which stocks a portfolio held mattered more for its terminal wealth than which utility
function weighted them, and the forecast band and the drawdown limit brought many re-plans forward of the horizon.
Compared with re-planning on the calendar alone, these early re-plans made the worst drawdown smaller in most
combinations of rules, on real prices and in the deepest simulated declines, but did not raise terminal wealth
consistently.
In the deepest simulated declines, every combination of selection and weighting rule ended ahead of both the
simulated market index and an equal-weight portfolio of all stocks, random selections included, and did so with the
early re-plans switched off as well. The equal-weight portfolio also ended ahead of the index, so part of the margin
came from the simulated market and the rest from parts of the system that all combinations shared, such as the
volatility cap, whose separate effect we did not test, rather than from the choice of stocks or the re-planning
triggers. The market hedge cost
wealth while markets rose and helped when declines persisted, and the synthetic futures favored score-based
selection partly because the generator preserved each stock's historical trend. More broadly, the results suggest
that portfolio systems are best evaluated as a whole, on synthetic as well as historical markets, and that a synthetic
test measures the value of stock selection fairly only when its generator does not build in the trend persistence
that selection rules exploit.
